\documentclass[11pt]{article}
\usepackage[margin=1in]{geometry}
\usepackage[T1]{fontenc}
\usepackage{lmodern}
\usepackage{microtype}
\usepackage{amsmath,amssymb}

\begin{document}

\section{Architecture}
\label{sec:architecture-classification}

\subsection{Architecture classification and justification}

The implementation in \texttt{agv\_hybrid.py} is most accurately described
as a \emph{hierarchical, decomposed planner with distributed floor-level
planning logic}, rather than as a fully decentralized multi-agent system.
The word ``decentralized'' on its own would overstate the current design:
the floor servers are separate planning components, but they are constructed
and orchestrated by one \texttt{run\_hybrid()} process, and the code does not
implement independent server processes or a peer-to-peer consensus protocol.

Several decisions remain centralized. The call to \texttt{algorithm\_1()}
assigns jobs and creates complete AGV routes before decomposition. The
central process splits those routes into legs, dispatches legs through a
single heap, calls \texttt{CentralServer.book\_lift()} to reserve every lift
track, and merges the floor-leg windows and lift intervals into final plans.
The central process also initiates final verification. These functions
provide global coordination and maintain the authoritative lift reservation
state.

Other decisions are spatially distributed across floor-specific planner
components. Each \texttt{FloorServer.plan\_leg()} plans a route leg on its own
floor, maintains that floor's local node reservations $W_n$ and edge
reservations $W'_e$, and invokes the route-window conflict-resolution
procedure using only the local reservation state. Conflicts within a floor
are therefore handled at the floor component, while inter-floor lift
contention is handled centrally. This division supports the terms
\emph{hierarchical} (central coordination above floor planners) and
\emph{decomposed} (route and reservation state partitioned by floor). It is
reasonable to describe the planning workload as distributed across floor
servers, but the current implementation is not fully decentralized because
assignment, lift arbitration, orchestration, and plan merging remain central.

The information flow can be summarized textually as follows:
\begin{verbatim}
                         jobs, full routes, leg dispatch
                     +------------------------------------+
                     |                                    v
               +-----------+                    +------------------+
               |  Central  | -- plan leg ------>| Floor server f=1 |
               |  server   | <-- local windows -| local W_n, W'_e  |
               |           |                    +------------------+
               | Algorithm 1| -- plan leg ------>| Floor server f=2 |
               | lift book | <-- local windows -| local W_n, W'_e  |
               | merge/check|                    +------------------+
               +-----------+ -- plan leg ------>| Floor server f=3 |
                    ^   |   <-- local windows -| local W_n, W'_e  |
                    |   +------------------------+------------------+
                    |          lift booking / handover
                    v
             +-------------+       execute assigned jobs/routes
             | AGV fleet   | <--------------------------------------+
             +-------------+
\end{verbatim}
In the implemented planner, the central-to-floor arrows are function calls
and returned data within one program, rather than network messages between
independent machines. AGVs receive completed schedules for execution; they
do not negotiate jobs or lift slots directly in the current implementation.
Relevant code is \texttt{run\_hybrid()} (\texttt{agv\_hybrid.py}, lines
106--165), \texttt{FloorServer.plan\_leg()} (lines 65--76), and
\texttt{CentralServer.book\_lift()} (lines 89--99).

\subsection{Design note: peer-to-peer lift reservation}

A minimal step toward stronger decentralization would remove the global
\texttt{CentralServer.lift} reservation table and make each lift reservation
a negotiation between the floor servers at that lift's two endpoints. The
job assignment and route decomposition could initially remain unchanged;
only lift arbitration would change. The departure-floor server would request
a slot from the arrival-floor server, and neither server would authorize
departure until both sides agree on the same reservation. This is a
design sketch only; it is not implemented in \texttt{agv\_hybrid.py}.

To prevent both endpoints from granting overlapping trips independently, the
protocol needs a consistent per-lift ordering and a commit handshake. A
reservation request can carry a unique trip identifier, lift identifier,
direction, AGV priority, earliest-ready time, traversal duration, and safety
separation. Both endpoint servers maintain a calendar of committed and
pending intervals. Contending requests are ordered deterministically, for
example by (earliest-ready time, AGV priority, request identifier). A
requester proposes the earliest interval consistent with that order; the
peer checks its calendar and returns an acknowledgment or a rejection with
the next feasible time. The pair commits the interval only after matching
acknowledgments for that request identifier. If a message is lost or an
endpoint fails before commit, the origin must not enter the lift; pending
holds expire only under a shared timeout/epoch rule understood by both
endpoints.

\paragraph{Pseudocode.}
\begin{verbatim}
Departure floor server, on AGV ready for lift e:
    create request id; send REQUEST(e, id, direction, ready, duration,
                                    priority, delta) to peer endpoint
    receive peer calendar state / ACK or NACK
    choose earliest candidate interval consistent with deterministic order
    send PROPOSE(e, id, interval) to peer
    if peer validates interval and returns ACK(e, id, interval):
        record COMMITTED interval; authorize AGV to depart
    else:
        defer AGV in departure lobby and retry at returned feasible time

Arrival floor server, on REQUEST or PROPOSE:
    serialize requests for e using the same deterministic ordering
    check committed and pending intervals against candidate, padded by delta
    ACK only if candidate is conflict-free and both endpoints can commit it
    on arrival / RELEASE(e, id), close the reservation at both endpoints
\end{verbatim}

This protocol introduces failure modes that the single central reservation
table avoids. Concurrent requests can cause \emph{double-booking} if endpoint
calendars diverge, acknowledgments are duplicated, or an old pending hold is
incorrectly expired. A failure after one side records a commit but before the
other side receives it can leave a ghost reservation or an unsafe departure
unless the commit protocol is idempotent and fail-safe. Delayed messages,
clock skew, inconsistent request ordering, or a server restart can also make
the two endpoints disagree about which slot is next. If reservations or
other resources are held while waiting for further negotiations, cyclic
waits can create \emph{deadlock}; repeated rejection of a low-priority
request can create starvation. A conservative protocol that stops when
coordination is uncertain improves safety but can reduce availability.

Validation should first compare peer reservations against the current
central booking policy on identical seeded workloads, using the independent
plan verifier extended to check lift intervals and endpoint agreement. Tests
should include simultaneous opposing-direction requests, several AGVs
requesting one lift, multiple lifts and floors, variable lift durations,
message delay/loss/duplication, endpoint restarts, and timeout/retry
behavior. Measurements should include verification failures, deadlocks,
double-bookings, reservation wait time, makespan, tardiness, throughput,
fairness across priorities, message count, and planning latency. Only after
these cases pass in simulation should the protocol be evaluated in a
hardware-in-the-loop setup and then on a physical lift-connected AGV system.

\end{document}
